\def\topic{}
\ifdefined\MyHeaderIncluded
\else
\def\qwertyflag{true}

\ifdefined\MyHeaderIncluded

\else
  \newcommand{\foobar}{FooBar}
  \def\MyHeaderIncluded{true}
% if is till the end of the file

\documentclass[10pt,a4paper]{article}
\usepackage[utf8]{inputenc}
%\usepackage[german]{babel}
%\usepackage[cyrillic]{babel}
\usepackage[T1]{fontenc}
%\usepackage[T2A]{fontenc}
%\usepackage[english, russian]{babel}
\usepackage{amsmath}
\usepackage{amsfonts}
\usepackage{amssymb}
\usepackage{amsthm}
\usepackage{mathtools} %coloneqq

\usepackage{bussproofs} %

%\renewcommand{\thesubsection}{\arabic{subsection}}
%\newtheorem{thm}{Theorem}[section]
%\renewcommand{\thethm}{\arabic{subsection}.\arabic{thm}}

\usepackage{thmtools}
\usepackage{hyperref}
\hypersetup{
    colorlinks=true,
    linkcolor=blue,
    filecolor=magenta,      
    urlcolor=cyan,
    pdftitle={Overleaf Example},
    %pdfpagemode=FullScreen,
    }
%\usepackage{stix-type1}%\turnediota
\theoremstyle{definition}

\newtheorem{constructor}{Constructor}[section]
\newtheorem{iaxiom}{Axiom}
\newtheorem{axiom}{Axiom}
\newtheorem*{axiom*}{Axiom}
\newtheorem{definition}{Definition}[section]
\newtheorem*{definition*}{Definition}
\newtheorem{remark}{Remark}[section]
\newtheorem*{remark*}{Remark}
\newtheorem{recall}{Recall}[section]
\newtheorem*{recall*}{Recall}

\newtheorem{lemma}{Lemma}[section]
\newtheorem{theorem}{Theorem}[section]
\newtheorem*{theorem*}{Theorem}
\newtheorem{note}{Note}[section]
\usepackage{etoolbox}

%\newcommand*{\@Dashv}[2]{%
%  \reflectbox{$\m@th#1#2$}%
%}
\newcommand*{\separation}[3]{% 1)binder, 2)set, 3)condition
  \{#1\in #2: #3\}
}

% https://www.freecodecamp.org/news/simple-regex-tricks-for-beginners-3acb3fa257cb/
%  \\mathcal\{T\} |---> \topology
%  \\mathcal\{([^\}]*)\}   |--->  $1
\newcommand{\AND}{\ \ \&\ \ }
\newcommand{\any}{\mathrm{\ any}}
\newcommand{\borel}{\mathcal{B}}
\newcommand{\can}{\mathrm{can}}
\newcommand{\canonize}{\mathrm{canonize}}
\newcommand{\closure}{\mbox{cl}}
\newcommand{\complex}{\mathbb{C}}
\newcommand{\converges}{\longrightarrow}
\newcommand{\countable}{\mbox{\ countable}}
\newcommand{\ldefined}{\mbox{\ defined}}
\newcommand{\dom}{\mbox{dom}}
\newcommand{\Dom}{\mbox{Dom}}
\newcommand{\Done}{\noindent \textbf{Done.}\\}
\newcommand{\lexists}{\mbox{\ exists}}
\newcommand{\lforall}{\mbox{\ forall}}
\newcommand{\lunique}{\mbox{\ unique}}
\newcommand{\family}{\mathrm{\ family}}
\newcommand{\Fm}{\mathrm{Fm}}
\newcommand{\FuncSymb}{\mathrm{FuncSymb}}
\newcommand{\function}{\mathrm{\ function}}
\newcommand{\functional}{\mbox{\ functional}}
\newcommand{\IFF}{\ \Longleftrightarrow\ }
\renewcommand{\iff}{\ \longleftrightarrow\ }
\newcommand{\IMPLIES}{\ \Longrightarrow\ }
\renewcommand{\implies}{\ \longrightarrow\ }
\newcommand{\Inf}{\mathrm{Inf}}
\newcommand{\injective}{\mbox{\ injective}}
\newcommand{\injection}{\mbox{\ injection}}
\newcommand{\interior}{\mbox{int}}
\newcommand{\intersection}{\bigcap}
\newcommand{\isempty}{\mathrm{\ empty}}
\newcommand{\LB}{\mathrm{LB}}
\newcommand{\nghbset}{\mathcal{N}}
\newcommand{\nonempty}{\mathrm{\ nonempty}}
\newcommand{\Ord}{\mathrm{Ord}}
\newcommand{\power}{\mathcal{P}}
\newcommand{\pr}{\mathrm{pr}}
\newcommand{\fst}{\mathrm{fst}}
\newcommand{\snd}{\mathrm{snd}}
\newcommand{\PredSymb}{\mathrm{PredSymb}}
\newcommand{\Proof}{\noindent \textbf{Proof.}\\}
\newcommand{\ran}{\mathrm{ran}}
\newcommand{\Ran}{\mathrm{Ran}}
\newcommand{\real}{\mathbb{R}}
\newcommand{\relation}{\mathrm{\ relation}}
\newcommand{\set}{\mathrm{\ set}}
\newcommand{\singleton}{\mathrm{\ singleton}}
\newcommand{\semantics}{\mathrm{semantics}}
\newcommand{\semialgebra}{\mathcal{S}}
\newcommand{\Sup}{\mathrm{Sup}}
\newcommand{\surjective}{\mbox{\ surjective\ on\ }}
\newcommand{\Term}{\mathrm{Term}}
\newcommand{\Type}{\mathrm{Type}}
\newcommand{\topology}{\mathcal{T}}
\newcommand{\total}{\mbox{\ total\ on\ }}
\newcommand{\transitive}{\mathrm{\ transitive}}
\newcommand{\union}{\bigcup}
\newcommand{\lTrue}{\mathrm{True}}
\newcommand{\lFalse}{\mathrm{False}}
\newcommand{\lUncertain}{\mathrm{Uncertain}}
\newcommand{\lif}{\mathrm{if}}
\newcommand{\lthen}{\mathrm{then}}
\newcommand{\lelse}{\mathrm{else}}
\newcommand{\lexception}{\mathrm{exception}}
\newcommand{\lundef}{\mathrm{undef}}
\newcommand{\univocal}{\mathrm{\ univocal}}
\newcommand{\univocalat}{\mathrm{\ univocal\ at\ }}
\newcommand{\UB}{\mathrm{UB}}
\newcommand{\Un}{\mathcal{U}}
\newcommand{\Var}{\mathrm{Var}}

\newcommand{\notvarphi}{\phi}
\renewcommand{\phi}{\varphi}

\newcommand{\comment}[1]{\quad\mbox{(#1)}}
% local commands
%\preto\section{\undef\blipo} %etoolbox
\preto\subsection{\undef\blipo} %etoolbox
% https://tex.stackexchange.com/questions/105684/limit-macro-newcommand-scope-to-section
\usepackage{graphicx}
\usepackage{authblk}

\setlength\parindent{0pt} % whole doc of \noindent
\usepackage{soul}
%\usepackage{mathgl}

\fi % guard
\ifdefined\topic
\else
\def\topic{No Titel}
\fi
\title{\topic}
\author[1]{Georgy Dunaev}
\begin{document}
\maketitle
\fi

\section{``Paradoxes'' of mathematics}
The arguments below separate inconsistent unrestricted comprehension from
valid class constructions. Each result states the logical or set-theoretic
principle it uses. In particular, the first two arguments are intuitionistic;
the regularity argument uses membership induction, the diagonal arguments
use separation, and the result about nonempty intersections uses classical
logic.

\subsection{There is no formula equivalent to its negation.}
In this subsection, $F$ is an arbitrary formula and negation is the special
case $C\implies\bot$. The arguments use implication introduction and
modus ponens, without the law of excluded middle.
\begin{lemma}[formula implies itself]
\label{lem:CtoC}
\[
(C\implies C)
\]
\Proof\ \\
\begin{enumerate}
\item \[(C\implies (C\implies C)) \tag{by ax.1}\]
\item \[(C\implies ((C\implies C)\implies C))\implies
(((C\implies (C\implies C))\implies (C\implies C)) )
\tag{by ax.2}\]
\item \[(C\implies ((C\implies C)\implies C)) \tag{by ax.1}\]
\item \[((C\implies (C\implies C))\implies (C\implies C)) \tag{by MP from steps 2 and 3}\]
\item \[(C\implies C)\tag{by MP from steps 1 and 4}\]
\end{enumerate}
\Done
\end{lemma}
\begin{note}[Implication and classical principles]
Via MP and the deduction theorem:
\[
C \implies ((C\implies F)\implies F)
\]
For comparison, Peirce's law is a classical principle and is not used in
the arguments here:
\[
(((P \implies Q) \implies P)\implies P)
\]
\end{note}

\begin{lemma}[Contraction of an implication]
\label{lem:CtoF}
\[
(C\implies (C\implies F))\implies (C\implies F)
\]
\Proof
Assume $(C\implies (C\implies F))$, and assume $C$.\\
By modus ponens from the assumptions, $(C\implies F)$.\\
Then $F$ is derived from $C$ and $(C\implies F)$ via modus ponens.
\Done
\end{lemma}

\begin{theorem}[Curry's paradox]
\label{Curry's Paradox}
\[
(C\iff (C\implies F))\implies F
\]
\Proof
Assume $(C\iff (C\implies F))$.\\
Therefore $(C\implies (C\implies F))$ and $((C\implies F)\implies C)$.\\
The formula $(C\implies F)$ holds by \myref{lem:CtoF} from $(C\implies (C\implies F))$.\\
%By axiom 2 and modus ponens from $ $,
Modus ponens applied to $((C\implies F)\implies C)$ and $(C\implies F)$
first gives $C$. A second application, to $C$ and $(C\implies F)$, gives $F$.
\Done
\end{theorem}

\subsection{There is no class containing all irreflexive classes}
\begin{theorem}[Russell's paradox]
\label{Russel's Paradox}
Both in classical and in intuitionistic logic, the following holds:
\[
\neg (\exists W \forall Q,\  Q\in W \iff Q\notin Q)
\]
\Proof
Assume that such a class $W$ exists and substitute $Q=W$. Put
$P\equiv(W\in W)$. The assumed equivalence gives $P\implies\neg P$
and $\neg P\implies P$. Assuming $P$ therefore gives both $P$ and
$\neg P$, so $\neg P$ follows by implication introduction. The second
implication now gives $P$, hence a contradiction. No case distinction
between $P$ and $\neg P$ is needed.
\Done
\end{theorem}
\begin{theorem}[The universal class is proper without regularity]
\label{thm:V_proper_separation}
Assume class comprehension and separation for subclasses of a set.
Then $\V$ is not a set. No foundation or choice axiom is needed.
\end{theorem}
\Proof
Form $R=\{x\in\V:x\notin x\}$. If $\V$ were a set, separation
would make $R$ a set, hence $R\in\V$. Its defining condition would
then give $R\in R\iff R\notin R$, contradicting the preceding
intuitionistic argument. Thus $\V$ is proper.
\Done

This conclusion is compatible with non-well-founded sets: it rules out
a universal set under separation, without ruling out individual membership
cycles. The next irreflexivity theorem uses the stronger hypothesis of
membership induction and should not be imported into settings lacking it.

%\subsection{Consequences of regularity axiom}
\subsection{Regularity implies that all classes are irreflexive}
\begin{theorem}[a class is not in itself]
With membership induction, the membership relation is irreflexive in both
classical and intuitionistic logic:
\[
\labeltagd{thm}{not in itself}
\forall Y,\ Y\notin Y
\]
\Proof
%\noindent %Fix Y.\\
Let $C\equiv \cc{W}{W\notin W}$.\\
Goal: to show that $(\power C\subseteq C)$.\\
Fix class $B$. Assume $B\in \power C$.\\
\[
(B\in \power C) \Longrightarrow (B\set)\land(B\subseteq C).
\]
Claim: $\forall b\in B,\ b\notin b$.\\ %Каждый его элемент не принадлежит себе.
$(b\in B)\LAND(B\subseteq C) \Longrightarrow (b\in C)$.\\
$(b\in C) \Longrightarrow (b\notin b)$.\\
Claim: $B\notin B$.\\
Assume opposite $B\in B$ and derive a contradiction.
\[
(B\in B) \Longrightarrow (B\notin B)
 \tag{by previous claim}
\]
%$(B\in B)\LAND(B\subseteq C) \Longrightarrow (B\in C)$.\\
%$(B\in C) \Longrightarrow (B\notin B)$.\\
$(B\in B)\LAND(B\notin B) \Longrightarrow \bot$.\\
Contradiction obtained, so the claim  $B\notin B$ is proven.\\
%Поэтому если бы B принадлежал себе(т.е. был своим элементом), то он бы не %принадлежал себе.
%Значит B не принадлежит себе.
$(B\set)\LAND(B\notin B) \Longrightarrow (B\in C)$.\\
%Значит B тоже элемент C(т.к. B — множество, которое не принадлежит себе).
%Доказали PC⊆C.
We are done proving the goal $\power C\subseteq C$.\\
\[
\power C\subseteq C \Longrightarrow V\subseteq C
\tag{by epsilon induction}
\]
%Значит из аксиомы следует, что A=V.
%Значит для любого множества W, W∉W.
%А если W не множество[т.е. собственный класс], то тогда он не принадлежит никакому классу по определению множества, как класса который чему--нибудь принадлежит.
%Значит любой класс не принадлежит сам себе.
It remains to cover arbitrary classes, including proper classes. Fix $Y$
and assume $Y\in Y$. Membership implies that $Y$ is a set, so $Y\in V$.
Since $V\subseteq C$, we obtain $Y\in C$ and hence $Y\notin Y$, a
contradiction. Thus $Y\notin Y$, without deciding whether $Y$ is a set.
\Done
\end{theorem}
\begin{theorem}[intersection set then nonempty]
In classical logic, if the intersection of a class $A$ is a set, then
$A$ is nonempty. Here the empty intersection is the universal class $\V$.
\[
\left(\left(\intersection A\right)\set\middle)
\implies
\middle(\exists w,\ w\in A\right)
\labeltagd{thm}{intersection set then nonempty}
\]
%\quad\mbox{(by lemma ... )}
\Proof
If $A$ has no member, extensionality gives $A=\emptyclass$, so
$\intersection A=\V$. Theorem~\ref{thm:V_proper_separation} says that
$\V$ is not a set, contradicting the hypothesis. Thus $A$ is not empty;
classical logic converts this into an existential witness $w\in A$.
The argument uses separation to establish properness of $\V$, but does
not use membership induction.
\Done
\end{theorem}

\subsection{The powerclass of a set \texorpdfstring{$A$}{A} is not a subclass of \texorpdfstring{$A$}{A}}
\begin{theorem}[powerset is not a subset]
In classical and even in intuitionistic logic, the following holds:
\[
\labeltagd{thm}{powerset not a subset}
\forall (Y\set),\ \neg(\power Y \subseteq Y)
\]
\end{theorem}
\begin{proof}
Use separation to form the set $D=\{x\in Y:x\notin x\}$. If
$\power Y\subseteq Y$, then $D\subseteq Y$ and sethood give
$D\in\power Y\subseteq Y$. Consequently
$D\in D\iff D\notin D$, a contradiction as in Russell's argument.
This proof requires neither excluded middle nor membership induction;
separation supplies the diagonal set.
\end{proof}

\begin{theorem}[Cantor's diagonal theorem]
For any set $A$, no function $f$ with domain $A$ can have every subset
of $A$ among its values. In particular, no function $A\to\power A$
is surjective. This uses separation but not choice or foundation.
\end{theorem}
\Proof
Define $D=\{x\in A:x\notin f(x)\}$ by separation. If $f(a)=D$
for some $a\in A$, then
\[
a\in D\iff a\notin f(a)\iff a\notin D,
\]
which is impossible intuitionistically. Thus the subset $D$ is missing
from the range. The proof does not need the powerclass to be a set;
the powerset axiom is needed only when treating that whole collection
as a set-valued codomain.
\end{proof}

\begin{remark}[Why the domain in Cantor's theorem is a set]
The book's powerclass $\power C$ collects the \emph{set} subclasses
of $C$. Consequently
\[
\power\V=\V:
\]
every set is a subclass of $\V$, and every member of a powerclass is
a set. The identity class function on $\V$ is therefore a surjection
onto $\power\V$. This does not contradict Cantor's theorem, whose
domain is a set. For that identity map the diagonal class is
$R=\{x\in\V:x\notin x\}$. It is a subclass of $\V$ but is proper:
if $R$ were a set, then $R\in\V$ and comprehension would give
$R\in R\iff R\notin R$. Thus $R$ is not a member of
$\power\V$ and cannot serve as a missing value within that codomain.
The collection of \emph{all} subclasses of $\V$, including proper
ones, is a different, conglomerate-level object.
\end{remark}

\subsection{The conglomerate of all classes is not a class}
\begin{theorem}[conglomerate of classes not class]
A conglomerate of all classes is not a class. (In the sense, that it ``does not exist as class''.)
\[
\neg (\exists W \forall Q,\ Q\in W)
\labeltagd{thm}{conglomerate of classes not class}
\]
\[
(\forall W \exists Q,\ Q\notin W)
\]
\Proof
For any class $W$, take $Q=\V$. Theorem~\ref{thm:V_proper_separation}
shows that $Q$ is a proper class; since membership in any class implies
sethood, $Q\notin W$. Thus the same witness works for every $W$.
This strengthens the earlier regularity-based argument: separation
suffices, even in the presence of non-well-founded sets.
\Done
\end{theorem}

\subsection{The class of all ordinals is not a set}

\begin{lemma}[On is not in On]
\[
\On\notin\On
\labeltagd{lem}{OnnotinOn}
\]
\end{lemma}%
\Proof
Assume $\On\in\On$.\\
Therefore $\in$ well--orders $\On$ as a well--founded strict linear order.\\
In particular strict order is irreflexive.\\
Under the assumption, $\On$ is itself an element of this strictly ordered
class, so irreflexivity gives $\On\notin\On$.\\
Contradiction with assumption.\\
Therefore $\On\notin\On$ is proven.
\Done

\begin{theorem}[the class of ordinal sets is not a set]
The class $\On$ of all ordinals is a proper class, i.e. it is not a set.
\[
\On\notin\V
\labeltagd{thm}{OnNotSet}
\]
%\label{thm:OnnotinV}
\Proof
Assume opposite: $\On\in\V$. To show: $\bot$.
\[
(\On\in\V)
\IMPLIES
(\On\set).
\]
By theorem \myref{thm:ordon}, $\Ord(\On)$. The rest is done as in the previous lemma:
\[
\Ord(\On)
\LAND
(\On\set)
\IMPLIES
(\On\in\On);
\]
\[
\Ord(\On)
\IMPLIES
(E\Irr\On);
\]
\[
(E\Irr\On)
\IMPLIES
(\On\notin\On);
\]
\[
(\On\in\On)
\LAND
(\On\notin\On)
\IMPLIES
\bot.
\]
\Done
\end{theorem}

\ifdefined\qwertyflag
\undef\qwertyflag
\end{document}
\fi
